\subsection{Prompts for Response Generation}
\textbf{TriviaQA}:
We use the exact prompt in~\cite{touvron2023llama} for TriviaQA, which is reproduced below:
{\scriptsize
\begin{verbatim}
Answer these questions:
Q: In Scotland a bothy/bothie is a?
A: House
Q: [Provided question]
A: 
\end{verbatim}
}

\textbf{Natural Questions} is a much harder dataset than TriviaQA, so we use the same 5-shot prompt version of the prompt in~\cite{touvron2023llama} (with 5 questions randomly picked from the training set).

%Q: who makes up the state council in russia
%A: governors and presidents
%Q: when does real time with bill maher come back
%A: November 9, 2018
%Q: where did the phrase american dream come from
%A: the mystique regarding frontier life
%Q: what do you call a group of eels
%A: bed
%Q: who wrote the score for mission impossible fallout
%A: Lorne Balfe

\textbf{CoQA}:
For CoQA, we use the code provided by~\cite{kuhn2023semantic} to prepare the data\footnote{https://github.com/lorenzkuhn/semantic\_uncertainty/blob/main/code/parse\_coqa.py}.
We provide the prompts below for convenience:
{\scriptsize
\begin{verbatim}
[The provided context paragraph]
[additional question-answer pairs]
Q: [Provided question]
A:
\end{verbatim}
}
where additional question-answer pairs are preceding turns of the conversation about
the paragraph consisting of questions and reference answers.

\subsection{Prompts for NLI similarity}
Similar to \citet{kuhn2023semantic}, we feed both the question and answer to DeBERTa-large\footnote{\url{https://huggingface.co/microsoft/deberta-large-mnli}} with the following prompt for any pair of answers to the same question in order to get the entailment/contradiction scores:
{\scriptsize
\begin{verbatim}
[question] [answer_1] [SEP] [question] [answer_2]
\end{verbatim}
}


\subsection{Automatic accuracy evaluation}%\label{appendix:check_judgement}
We use \texttt{gpt-3.5-turbo} to evaluate the similarity between each response and the reference answer.
The prompts for \datasetcoqa, \datasettrivia and \datasetnqopen are shown below. The few-shot examples are chosen from the training split.

\textbf{CoQA}:
%
{\scriptsize
%\begin{lstlisting}[breaklines]
\begin{verbatim}
Rate the level of consistency between the answer to the question and the reference answer, from 0 to 100.
Question: When was the Vat formally opened?
Reference: It was formally established in 1475
Answer: In 1475
Rating: 100.

Question: what is the library for?
Reference: research
Answer: tourism
Rating: 0.

Question: [question]
Reference: [reference answer]
Answer: [generated response]
Rating:
\end{verbatim}
%\end{lstlisting}
}

\textbf{TriviaQA}:
%\begin{lstlisting}[breaklines]
{\scriptsize
\begin{verbatim}
Rate the level of consistency between the answer to the question and the reference answer, from 0 to 100.
Question: In Scotland a bothy/bothie is a?
Reference: House
Answer: House
Rating: 100.

Question: Where in England was Dame Judi Dench born?
Reference: York
Answer: London
Rating: 0.

Question: [question]
Reference: [reference answer]
Answer: [generated response]
Rating:
\end{verbatim}
}
%\end{lstlisting}

\textbf{Natural Questions}:

%\begin{lstlisting}[breaklines]
{\scriptsize
\begin{verbatim}
Rate the level of consistency between the answer to the question and the reference answer, from 0 to 100.
Question: who makes up the state council in russia
Reference: governors and presidents
Answer: governors and presidents
Rating: 100.

Question: when does real time with bill maher come back
Reference: November 9, 2018
Answer: September 8, 2000
Rating: 0.

Question: [question]
Reference: [reference answer]
Answer: [generated response]
Rating:
\end{verbatim}
}
%\end{lstlisting}

99.34\% of the judgements by GPT can be parsed as an integer between 0 and 100 in the first attempt with a simple \texttt{str.split} in Python. 
We skipped the remaining samples for the experiments in this paper, but if necessary one could devise a better parsing algorithm as well as improved prompts to elicit judgement with higher coverage.



\textbf{Verifying the correctness of GPT evaluations}
We sample 33 samples per dataset and model (OPT, LLaMA, GPT) and perform a human evaluation of the quality of the GPT judgement. 
That is, we compare the human judgement on whether a generated response is correct or not with the GPT's judgement, like in \citet{kuhn2023semantic}.
In \cref{table:appendix:gpt_eval}, we show the breakdown of the accuracy by datasets.

\begin{table}[ht]

\caption{
The accuracy of the GPT evaluation on the correctness of the responses, measured by the authors. 
For example, for \datasettrivia, the human evaluations and the GPT evaluations are aligned on 97 out of the 99 question/answer pairs, leading to an accuracy of 98.0. 
}
\label{table:appendix:gpt_eval}
%\begin{center}
\centering
\begin{small}
% \scalebox{0.8}{
   \resizebox{0.6\columnwidth}{!}{
\begin{tabular}{l|ccc}
\toprule
 & \datasetcoqa & \datasettrivia & \datasetnqopen\\
\midrule
Accuracy of GPT evaluation & 90.9 & 98.0 & 95.9\\
\bottomrule
\end{tabular}
}
\end{small}
%\end{center}
\end{table}

We also include a few typical examples where the GPT evaluation is more reliable than the rougeL evaluation used in \citet{kuhn2023semantic}
\begin{itemize}
    \item ``Q: Were they nice to the other critters? A: no''. GPT correctly identifies that ``They were mean to the other animals'' is a correct answer, but the rougeL metric does not capture this logic.
    
    \item ``Q: Who created the immunity plan? A: the Gulf Cooperation Council''. 
    GPT correctly identities that ``GCC'' is a correct answer (an abbreviation of the Gulf Coorperation Councol'', but rougeL cannot identify this mechanically).

    \item ``Q: when was the last bear killed in the uk A: c. 1000 AD''.
    GPT classifies the response ``The last bear in the UK was killed over a thousand years ago, in the 9th century, so there is no specific date available'' as correct, but rougeL fails to do so.
    
\end{itemize}
The GPT evaluation also introduces its own problems, though. 
For example, when it is used to evaluate its own response to the question ``when did it become law to stand for the national anthem'' with reference answer ``6/22/1942'', it rates ``There is no federal law that mandates standing for the national anthem in the United States'' as 100.
However, we believe it is using its own knowledge here and largely ignores the reference answer. 
Such problems could potentially be improved with better prompts.


\subsection{Hyper-parameters for Uncertainty/Confidence Measures}
For all $U$ and $C$ measures involving $\affinityEntail$ and $\affinityContradict$, we need to choose a temperature for the NLI model.
The temperature is chosen from 0.1, 0.25, 0.5, 1, 3, 5, and 7.
For $U_\methodEcc$ and $C_\methodEcc$, we also need to choose a cutoff for eigenvalues.
For simplicity we use the same threshold for each experiment/dataset, and the threshold is chosen from 0.4, 0.5, 0.6, 0.7, 0.8, and 0.9.